\documentclass[book.tex]{subfiles}
\begin{document}

\chapter{Adsorption}
\label{chap:adsorption}
\noindent\rule{11cm}{0.4pt} \\
\textbf{Prerequisites:} \autoref{chap:quantum_mechanics}, \autoref{chap:stat_thermo} \\
\textbf{Difficulty Level:} ** \\
\noindent\rule{11cm}{0.4pt}
\vspace{5mm}

Adsorption occurs when atoms, ions, or molecules from a gas, liquid, or dissolved solid become bound to a surface due to favorable interactions (\textit{e.g.} Van der Waal's forces), creating a film of the adsorbate on the surface of the adsorbent. The term should not be confused with \textit{absorption}, in which a fluid (the absorbate) permeates or is dissolved by a liquid or solid (the absorbent). Adsorption considers a surface-based process, while the later involves the whole volume of the material.
The thermodynamic description of adsorption varies according to the characteristics of the adsorbent. In this chapter we begin with a study of non-interacting adsorption, ignoring the energy of interaction between the adsorbed components. We consider two cases: mobile and localized adsorption. We end by discuss phenomena arising from the interaction between adsorbed molecules 


\begin{marginfigure}
% \includegraphics[width=\linewidth]{figures/Physics/statistical_mechanics/adsorption/Fig1.png}
\includegraphics[width=\linewidth]{figures/Physics/statistical_mechanics/adsorption/CO_adsorbing_on_Ni.png}
\caption{Images from a molecular simulation of $CO$ adsorbing on a $Ni(110)$ surface (see \textit{http://www.fhi-berlin.mpg.de/\textasciitilde hermann/hermann/Ba lsacPictures/CONi110.gif}).}
\label{fig1}
\end{marginfigure}

\section{Non-interacting adsorption}

We first address adsorption on a lattice where there is negligible interaction between the adsorbed molecules. 

Consider the following two scenarios for the surface adsorption (Fig.~\ref{fig2}):

\begin{itemize}
\item \textit{Localized adsorption} occurs when the binding potential exhibits large spatial variations causing the adsorbate to be highly localized.
\item \textit{Mobile adsorption} occurs when the surface potential is nearly homogeneous and the adsorbate is highly mobile.
\end{itemize}

\begin{marginfigure}
% \includegraphics[width=\linewidth]{figures/Physics/statistical_mechanics/adsorption/Fig2.png}
\includegraphics[width=\linewidth]{figures/Physics/statistical_mechanics/adsorption/2_adsorption (1).png}
\caption{Schematic of two scenarios of adsorption.}
\label{fig2}
\end{marginfigure}

In localized adsorption, each adsorbed molecule is trapped in a potential well located at a specific lattice site. To a first approximation, we can expand the potential in the $x$, $y$, and $z$ directions to quadratic order in the displacement. Thus, the single-molecule partition function of an adsorbed molecule
is:
\begin{align}\label{Eq1}
q=q_xq_yq_ze^{-\beta V_0} \hspace{0.4cm} \text{with} \hspace{0.4cm} q_i = \frac{e^{-\theta _{vib,i}/(2T)}}{1-e^{-\theta _{vib,i}/(2T)}}
\end{align}

where $q_i(i =x,y,z)$ is a 1-dimensional harmonic oscillator partition function, \begin{align}
\theta_{vib,i} = \frac{\hslash \nu_i}{k_B}
\end{align}
is the vibrational temperature in the $i$ direction, and $V_0$ is the potential at the well bottom.

Since each adsorbed molecule occupies a lattice site with a distinct spatial location, they are distinguishable. For a lattice with $M$ sites containing N adsorbed molecules, the degeneracy of configurations is 
\begin{align}
\Omega = \frac{M!}{N!(M-N)!}
\end{align}
and the canonical partition function for localized adsorption is:
\begin{align}
Q_{localized}(N,M,T) =\frac{M!}{N!(M-N)!}q^N
\end{align}

The chemical potential of the adsorbate is given by:
\begin{align}
\frac{\mu}{k_BT} &= -\left( \frac{\partial \log Q_{localized}}{\partial N}\right)_{M,T} \\
&= -\log q + \log N -\log(M-N) \\
&= \log \left[ \frac{\theta}{(1-\theta)q} \right] \hspace{0.3cm} \text{where} \hspace{0.3cm} \theta = \frac{N}{M}
\end{align}

Assuming adsorption from the gas phase, equilibrium dictates equality of the chemical potentials between the adsorbate on the surface and in the gas phase. For an ideal gas, the chemical potential can be written as 
\begin{align}
\frac{\mu}{k_BT} &= \frac{\mu_0(T)}{k_BT} + \log p
\end{align}
This leads to the \textit{Langmuir adsorption isotherm}:
\begin{align}\label{Eq4}
\theta(p,T) = \frac{\chi(T)p}{1+\chi(T)p} \hspace{0.3cm} \text{where} \hspace{0.3cm} \chi(T) = q(T)e^{\beta\mu_0(T)} 
\end{align}

\noindent that describes the surface coverage $\theta$ versus the temperature and pressure of the gas-phase adsorbent for localized adsorption on a solid adsorbent. The coverage tends to unity when $p \rightarrow \infty$ (total site coverage).

In the case of mobile adsorption, the single-molecule partition function will have a $z$ direction contribution and an $xy$ contribution that reflects surface mobility. This is constructed by assuming a 2D ideal gas with a separate z
component partition function, \textit{i.e.}:

\begin{align}\label{Eq5}
q = \frac{A}{\Lambda^2}q_z^{-\beta V_0} \hspace{0.3cm} \text{where} \hspace{0.3cm} \Lambda = \sqrt{\frac{h^2}{2\pi m k_BT}}
\end{align}

\noindent and $A$ is the surface area of the adsorbent.

The mobile adsorbed molecules are indistinguishable, since they are not restricted to a localized adsorption site. The canonical partition function is therefore given by:
\begin{align}\label{Eq6}
Q_{mobile}(N,A,T) = \frac{1}{N!}q^N
\end{align}
This gives the chemical potential:

\begin{align}\label{Eq7}
\frac{\mu}{k_BT} = -\left( \frac{\partial \log Q_{mobile}}{\partial N}\right)_{M,T} = -\log q + \log N 
\end{align}

For adsorption from an ideal gas, the surface coverage is found to be:

\begin{align}\label{Eq8}
\log\left(\frac{N}{q} \right) = \log\left(pe^{\beta \mu_0} \right) \rightarrow \frac{N}{A} = \chi'(T)p
\end{align}

\noindent where 
\begin{align}
\chi'(T) &= \frac{1}{\Lambda^2}q_ze^{-\beta V_0}e^{-\beta \mu_0}
\end{align}

All things the same, this leads to $N_{mobile} > N_{localized}$ at all pressures $p>0$ (Fig.~\ref{fig3}).

\begin{marginfigure}
\includegraphics[width=\linewidth]{figures/Physics/statistical_mechanics/adsorption/Fig3.png}
\caption{Plot of the surface coverage for localized adsorption ($\theta$ versus
$\chi,p$, solid curve) and mobile adsorption ($N=A$ versus $\chi'p$, dashed curve)}
\label{fig3}
\end{marginfigure}

\section{Multi-component adsorption}

This is easily extended to a system with $s$ component (localized adsorption).

Define $q_i$ as the single-molecule partition function of the s component, and $N_i$ as the number of adsorbed molecules of the $i$th adsorbate (note, $N_{tot} = \sum_{i=1}^s N_i\leq N$). The degeneracy of arranging $N_1,N_2,...,N_s$ adsorbed molecules on $M$ sites is 
\begin{align}\Omega = \frac{M!}{N_1!N_2!...N_s!(M - N_{tot})!}\end{align} 
The canonical partition function is given by:

\begin{align}\label{Eq9}
Q(N,M,T) = \frac{M!}{N_1!N_2!...N_s!(M - N_{tot})!} q_1^{N_1}q_2^{N_2}...q_s^{N_s}
\end{align}

The chemical potential of the $i$th species is:

\begin{align}
\frac{\mu_i}{k_BT} &= -\left( \frac{\partial \log Q}{\partial N_i}\right)_{N_{j \neq i},M,T} \notag\\
&= -\log q -\log(M-N_{tot}) + log N_i \notag\\
&= \log \left[ \frac{N_i}{(M-N_{tot})q_i} \right] \notag\\
&= \log \left[ \frac{\theta_i}{(1-\sum_{j=1}^s \theta_j)q_i} \right] \notag
\end{align}

For equilibrium with an ideal gas phase, the chemical potential of the $i$th species is written as 
\begin{align}
\frac{\mu_i}{k_BT} &= \frac{\mu_{i,0}(T)}{k_BT} + \log p_i
\end{align}
where $p_i = py_i$ is the partial pressure of the $i$th species ($y_i$ is the mole fraction in the gas phase).

Defining 
\begin{align}
\chi_i(T) &= q_ie^{\beta \mu_{i,0}}
\end{align}, 
we have
\begin{align}
p_i\chi_i = \frac{\theta_i}{1 - \sum_{j=1}^s \theta_j} \hspace{0.3cm} \text{where} \hspace{0.3cm} \theta_i = \frac{p_i\chi_i}{1 + \sum_{j=1}^s p_j\chi_j}
\end{align} 

The two extremes of the multi-component isotherm are as follows: For $p_i\chi_i \rightarrow 0$, the coverage behaves as 
\begin{align}
\theta_i \approx p_i\chi_i
\end{align}
which is the Henry's law regime.
For $p_i\chi_i \rightarrow \infty $ while $p_j\chi_j$ remains finite ($i \neq j$), the coverage saturates to 
\begin{align}\theta_i \rightarrow 1
\end{align}
Which is the complete coverage regime.

Therefore, binding affinity for a particular species can be outcompeted by species with lower affinity by raising the partial pressure of the other components

\section{Adsorption with interactions between adsorbed molecules}
Our previous analysis neglects any interactions between adsorbed molecules, which is an adequate description for low surface coverage and lattice structure with large spacing. However, surface interactions leads to correlated surface coverage,
analogous to the spatial correlations that we studied with the Ising model.

Let's adopt a simple model for the adsorption energy on a lattice:

\begin{align}\label{Eq11}
\beta E = f_{bind}\sum_{i=1}^M n_i - K \sum_{\langle i j \rangle} n_in_j
\end{align}

\noindent where $n_i$ is the site occupancy ($n_i = 0$ if unoccupied and $n_i = 1$ is occupied), $K$ is a coupling constant to describe interactions between neighboring adsorbed molecules, and $f_{bind}$ is the free energy of binding the adsorbed molecule (note, $f_{bind} = -k_BT \log q$).

The grand canonical partition function is written as:
\begin{align}
\Xi = \sum_{n_1,n_2,...n_M = 0,1} \exp \left( \beta K \sum_{\langle i j \rangle} n_in_j - \beta f_{bind}\sum_{i=1}^M n_i + \beta \mu\sum_{i=1}^M n_i \right)
\end{align}

Rather than analyze this further, let's convert this into a form that demonstrates that this is exactly the Ising model that we have already studied extensively. We relate the occupancy number $n_i$ to the spin $s_i$ using $n_i = \frac{s_i+1}{2}$. which gives $n_i=1$ when $s_i=1$ and $n_0=0$ when $s_i=-1$. With this we write:
\begin{align}
&\beta K \sum_{\langle i j \rangle} n_in_j - \beta f_{bind}\sum_{i=1}^M n_i + \beta \mu\sum_{i=1}^M n_i = \\
&\frac{\beta K}{4} \sum_{\langle i j \rangle} s_is_j + \beta \left(Kz + \frac{\mu}{2} - \frac{f_{bind}}{2}\right)\sum_{i=1}^M s_i + \beta M \left(\frac{Kz}{4}+\frac{\mu}{2} - \frac{f_{bind}}{2}\right) 
\end{align}

\noindent where $z$ is the coordination number of the lattice giving the number of near neighbors divided by 2 to avoid double counting (\textit{e.g.} $z = 2$ for a 2D square lattice).

Therefore, we can exactly map our current problem onto the Ising model with the parameters $J = K/4$ and $h = Kz+\mu/2- f_{bind}/2$, thus leveraging the work that we have already done.

Owing to the exact mapping between our current problem and the Ising model, identical physical phenomena arise in both systems
\begin{itemize}
\item Adsorption with interactions exhibits a phase transition under some conditions (sufficiently large coupling constant $K$ and lattice that is not 1 dimensional).
\item Adsorption exhibits hysteresis in the surface coverage versus gas pressure associated with this phase transition.
\item The surface coverage is not completely uncorrelated; rather, the adsorbed molecules exhibit correlated fluctuations due to favorable interactions.
\item The critical point in the adsorption isotherm is associated with a divergence in the correlation length of these spatial fluctuations.
\end{itemize}


%\section{Ising model approach to adsoprtion reactions in \\PEMFCs}
%Adsorption of water and its intermediate species like -OH, -O$_2$ on the (111) termination of Pt metal catalyst is the most studied problem in the area of PEMFCs. The (111) termination of Pt can be understood as a 2D triangular lattice where each site can be either empty, or is occupied by wither an adsorbed OH, or H$_2$O as shown in Fig.~\ref{ads}.

%It is very difficult to solve a 2D Ising model exactly and therefore, we will look at this problem from a \textit{mean-field} approach of the form:
%\begin{align}
%     \beta E = -h\sum_{i=1}^{N}\sigma_{i} - Jzm\sum_{i=1}^{N}\sigma_{i} \notag
%\end{align}

%\begin{marginfigure}
%\includegraphics[width=\linewidth]{figures/part2a/statistical_mechanics/adsorption/ads.png}
%\caption{Snapshot of the adsorbate layer on Pt(111) formed from the simulated voltammogram at 0.85 V RHE. The Pt atoms are green, O are red, and H are white. From Viswanathan et al. J. Phys. Chem. C  2012, 116, 4698-4704.}
%\label{ads}
%\end{marginfigure}

%where the symbols have their usual meanings. We consider three possible values of spin depending on the site $i$ being either empty ($\sigma_{i}^{1}$) or being occupied by an adsorbed OH ($\sigma_{i}^{2}$) or an adsorbed H$_2$O ($\sigma_{i}^{3}$). For an empty site the spin variables will have the value $\sigma_{i}^{1} = 1$, $\sigma_{i}^{2} = \sigma_{i}^{3} = 0$. Likewise, for a OH adsorbed site the spin variables will have the value $\sigma_{i}^{2} = 1$, $\sigma_{i}^{1} = \sigma_{i}^{3} = 0$ and for a H$_2$O adsorbed site the spin variables will have the value  $\sigma_{i}^{3} = 1$, $\sigma_{i}^{1} = \sigma_{i}^{2} = 0$. The empty sites do not respond to any external field and have h$^{1}$ = 0 eV, h$^{2}$ = -0.6 eV and h$^{3}$ = 0.2 eV.

%Considering only the nearest neighbor interactions, J$^{kl}$ denotes the coupling constant for interaction between two neighbor sites where one of the sites has occupation $\sigma^{k}$ and the other has occupation $\sigma^{l}$. The coupling between an empty site and any other site has no contribution, therefore J$^{11}$ = J$^{12}$ = J$^{13}$ = J$^{21}$ = J$^{31}$ = 0 eV. The OH -- OH interaction is known to be unfavorable and its coupling constant is known to be  J$^{22}$ = -0.3 eV. Likewise the H$_2$O -- H$_2$O interaction is also unfavorable and its coupling constant is given as J$^{33}$ = -0.3 eV. The interaction between OH and H$_2$O, however, is slightly favorable and has a coupling constant J$^{23}$ = J$^{32}$ = 0.1 eV. Using a value of kT = 0.0259 eV, solve the following exercises.

\section{Exercises}

\begin{enumerate}
\item Within the mean-field approach, write the expression for the Hamiltonian H = $\beta E$ for the total system of N sites on a 2D triangular lattice in terms of the various spin variables, external fields and coupling constants (without substituting actual values) and average occupation, where the average occupation of kind 'k' can be denoted as $\langle \sigma^{k}\rangle $. 
\item Express the partition function $Q$ for this system of N sites using the defined Hamiltonian in terms of the the various spin variables, external fields, coupling constants and average occupations. 
\item What is the expression for each of the average magnetization $\langle \sigma^{k}\rangle $ in terms of the various spin variables, external fields, coupling constants and themselves. 
\item Is this system capable of phase transition?
\end{enumerate}

\end{document}